\subsection{映入度量空间的映射空间}

设 X 是一个集合，Y 是一致空间，\((f_{t})_{t\in I}\) 是定义 Y 的一致结构的伪度量族（第 IX 章，§ 1，第 4 号），S 是 X 的子集所成之集。对每个 \(t\in I\) 、每个集 \(A\in S\) 以及每对 X 到 Y 中的映射 \((u,v)\) ，记

\[
g _ {t, \mathrm{A}} (u, v) = \sup _ {x \in \mathrm{A}} f _ {t} (u (x), v (x));
\]

由此立即得出，\(g_{t,\mathbf{A}}\) 是 \(\mathcal{F}(\mathbf{X};\mathbf{Y})\) 上的一个伪度量，且伪度量族 \((g_{t,\mathbf{A}})_{t\in \mathbf{I},\mathbf{A}\in \mathfrak{S}}\) 定义了 \(\mathfrak{S}\) -收敛的一致结构（在 \(\mathcal{F}(\mathbf{X};\mathbf{Y})\) 上）。特别地：

命题 1. 若 Y 是可度量化的一致空间，则 \(\mathcal{F}(X;Y)\) 上的一致收敛的一致结构是可度量化的。

因为若 d 是 Y 上与其一致结构相容的一个度量，则 \(\mathcal{F}(X;Y)\) 上的一致收敛结构由单个伪度量

\[
\delta (u, v) = \sup _ {x \in \mathbf {X}} d (u (x), v (x));
\]

一般说来这个伪度量不是有限的，但它等价于一个有限伪度量（第 IX 章，§ 1，第 2 号）；由于一致收敛的一致结构是豪斯多夫的（§ 1，第 2 号，命题 1），故它是可度量化的。

推论. 设 X 是拓扑空间，Y 是可度量化的一致空间。假设存在 X 的紧子集的一个序列 \((\mathbf{K}_{n})\) ，使得 X 的每个紧子集都含于某个 \(K_{n}\) 。则 \(\mathcal{F}(\mathbf{X};\mathbf{Y})\) 上的紧收敛的一致结构是可度量化的。

由于诸 \(\mathbf{K}_n\) 覆盖 \(\mathbf{X}\) ，\(\mathcal{F}_c(\mathbf{X};\mathbf{Y})\) 同构于 \(\prod_{n}\mathcal{F}_{u}(\mathbf{K}_{n};\mathbf{Y})\) 的一个一致子空间（ \(\S 1\) ，第 2 号，注 3），因此该推论

由命题 1 得出（第 IX 章，§ 2，第 4 号，定理 1 的推论 2）。

注意，当 X 局部紧且 σ-紧时（第 I 章，§ 9，第 9 号，命题 15 的推论 1），这个推论特别适用。

现在设 Y 是度量空间，d 是它的度量。若 X 是任一集合且 S 是 X 的子集所成的任一集，我们将用 \(\mathcal{B}_{\mathfrak{S}}(X;Y)\) 表示所有映射 \(u:X\to Y\) （使得 \(u(A)\) 对每个 \(A\in S\) 有界）所成之集。除非明确说明相反的情形，我们把 \(\mathcal{B}_{\mathfrak{S}}(X;Y)\) 看作赋有 S-收敛的一致结构，它由 \(\mathcal{B}_{\mathfrak{S}}(X;Y)\) 上的下列伪度量族定义：

\[
d _ {\mathbf {A}} (u, v) = \sup _ {x \in \mathbf {A}} d (u (x), v (x))\tag{A ∈ S}
\]

由假设这些伪度量都是有限的。当 \(\mathfrak{S} = \{\mathbf{X}\}\) 时，我们用 \(\mathcal{B}(\mathbf{X};\mathbf{Y})\) 代替 \(\mathcal{B}_{\mathfrak{S}}(\mathbf{X};\mathbf{Y})\) 。映射 \(u:\mathbf{X}\to \mathbf{Y}\) 称为有界的，如果它属于 \(\mathcal{B}(\mathbf{X};\mathbf{Y})\) ，即如果 \(u(\mathbf{X})\) 是 \(\mathbf{Y}\) 的一个有界子集。

命题 2. 设 X 是一个集合，Y 是度量空间。有界映射所成之集 \(\mathcal{B}(X;Y)\) 在空间 \(\mathcal{F}_{u}(X;Y)\) 中既是开的又是闭的。

若 \(u\) 有界，则每个映射 \(v: \mathbf{X} \to \mathbf{Y}\) （满足对一切 \(x \in \mathbf{X}\) 有 \(d(u(x), v(x)) \leqslant 1\) ）都是有界的，因为

\[
d (v (x), v (x _ {0})) \leqslant d (u (x), u (x _ {0})) + 2;
\]

因此 \(\mathcal{B}(\mathbf{X};\mathbf{Y})\) 是开的。另一方面，若 \(u\) 属于 \(\mathcal{B}(\mathbf{X};\mathbf{Y})\) 在 \(\mathcal{F}_u(\mathbf{X};\mathbf{Y})\) 中的闭包，则存在映射 \(u_0\in \mathcal{B}(\mathbf{X};\mathbf{Y})\) ，使得 \(d(u(x),u_0(x))\leqslant 1\) 对一切 \(x\in \mathbf{X}\) 成立；因此 \(u\) 有界。

推论 1. 设 X 是一个集合，Y 是度量空间。则 \(\mathcal{B}_{\mathfrak{S}}(X;Y)\) 在 \(\mathcal{F}_{\mathfrak{S}}(X;Y)\) 中是闭的。特别地，若 Y 完备，则 \(\mathcal{B}_{\mathfrak{S}}(X;Y)\) 关于 \(\mathfrak{S}\) -收敛的一致结构是完备的。

因为 \(\mathcal{B}_{\mathfrak{S}}(\mathbf{X};\mathbf{Y})\) 是子集 \(\prod_{\mathbf{A}\in \mathfrak{S}}\mathcal{B}(\mathbf{A};\mathbf{Y})\) （积 \(\prod_{\mathbf{A}\in \mathfrak{S}}\mathcal{F}_u(\mathbf{X};\mathbf{Y})\) 的子集）在典型映射（由 \(\mathcal{F}_{\mathfrak{S}}(\mathbf{X};\mathbf{Y})\) 到 \(\prod_{\mathbf{A}\in \mathfrak{S}}\mathcal{F}_u(\mathbf{A};\mathbf{Y})\) 中）下的原像；于是第一个断言由 § 1，第 2 号，注 3 得出，而第二个断言由第一个推出，只要注意到 § 1，第 5 号的定理 1。

推论 2. 设 X 是拓扑空间，Y 是度量空间。则 X 到 Y 中的一切有界连续映射所成的空间在 \(\mathcal{C}_u(X;Y)\) 中既是开的又是闭的；若 Y 完备，则它是完备的。

所考虑的空间是 \(\mathcal{B}(\mathbf{X};\mathbf{Y})\cap \mathcal{C}_u(\mathbf{X};\mathbf{Y})\) ；第一个断言由命题 2 得出；第二个断言由第一个推出（§ 1，第 6 号，定理 2 的推论 1）。

